\chapter{Array e puntatori}
\label{chap:puntatori}

I programmi lavorano spesso su molti dati dello stesso tipo: i voti di una
classe, le misure di un esperimento, le parole di un testo. Gli
\textbf{array} sono il tipo di dato che li raccoglie in variabili
\textbf{indicizzate}. In C gli array sono strettamente legati ai
\textbf{puntatori}: il nome di un array è trattato come un puntatore
costante, i puntatori si possono indicizzare come gli array, e il C ha
un'\textbf{aritmetica dei puntatori}. I puntatori servono anche a ottenere
l'effetto del passaggio “per riferimento” dei parametri, e a chiedere memoria
mentre il programma è in esecuzione. Le stringhe, che in C sono array di
caratteri, e gli array a più dimensioni sono l'argomento del
capitolo~\ref{chap:stringhe}.

\section{Gli array}

\subsection{Tante variabili, un solo nome}

Un programma che gestisce dei voti potrebbe dichiarare
\code{int grade0, grade1, grade2;}. Ma se le variabili necessarie sono
cento, o mille, dichiararle tutte una per una diventa impossibile. Una
variabile \textbf{indicizzata} è molto più comoda:

\begin{lstlisting}[style=c]
int grade[3];                        /* un array di tre interi */

grade[0] = grade[1] = grade[2] = 0;
\end{lstlisting}

\code{grade} è un \term{array} di tre interi, che si chiamano
\code{grade[0]}, \code{grade[1]} e \code{grade[2]}: un solo nome, e un
\textbf{indice} fra parentesi quadre che seleziona l'elemento. Il primo
indice è \textbf{sempre 0}. In memoria gli elementi occupano celle
\textbf{consecutive} (figura~\ref{fig:tre-variabili}).

\begin{figure}[H]
  \centering
  \begin{tikzpicture}[font=\Lato\scriptsize]
    \node[font=\Lato\scriptsize\bfseries] at (1.4,1.0) {tre variabili separate};
    \variabile{g0}{(1.4,0.3)}{grade0}{0}
    \variabile{g1}{(1.4,-0.3)}{grade1}{0}
    \variabile{g2}{(1.4,-0.9)}{grade2}{0}
    \node[font=\Lato\scriptsize\bfseries] at (6.6,1.0) {un array: \code{int grade[3];}};
    \node[mnome, anchor=east] at (5.2,-0.3) {grade};
    \celle[minimum width=1.0cm]{g}{(5.3,-0.3)}{0,0,0}
    \node[mnota, anchor=west, align=left] at (8.6,-0.3) {celle consecutive;\\sotto, l'indice di ciascuna};
  \end{tikzpicture}
  \caption{Tre variabili indipendenti e un array di tre elementi.}
  \label{fig:tre-variabili}
\end{figure}

\subsection{La dichiarazione}

La dichiarazione di un array a una dimensione indica il \textbf{tipo} degli
elementi, l'\textbf{identificatore} e, fra parentesi quadre, un'espressione
\textbf{costante} intera: la dimensione. La dichiarazione
\code{int a[size];} riserva lo spazio per gli elementi da \code{a[0]} a
\code{a[size-1]}. È buona norma definire la dimensione come costante
simbolica, così che per cambiarla basti modificare una riga:

\begin{affianco}[0.5][c]
\begin{lstlisting}[style=c, name=somma.c]
#include <stdio.h>

#define N 4

int main(void)
{
  int a[N] = {3, 7, 1, 9}, i;
  int sum = 0;

  for (i=0; i<N; i++)
    sum += a[i];
  printf("Sum is %d \n", sum);
  return 0;
}
\end{lstlisting}
\risultato
\begin{console}
$ ./somma
Sum is 20
\end{console}
\vspace{2pt}
\wtnota{Il ciclo \code{for} con \code{i} da 0 a \code{N-1} è il modo tipico
di visitare tutti gli elementi.}
\end{affianco}

Come per ogni variabile, ogni elemento occupa i byte del suo tipo: gli
elementi di un array di \code{int} stanno a 4 byte di distanza l'uno
dall'altro (figura~\ref{fig:array-memoria}).

\begin{figure}[H]
  \centering
  \begin{tikzpicture}[font=\Lato\scriptsize]
    \memoria[minimum width=1.6cm]{m}{(2,0)}{1000/3/a[0], 1004/7/a[1], 1008/1/a[2], 1012/9/a[3]}
    \draw[wtbrace] ([xshift=1.9cm]m-1.north east) -- ([xshift=1.9cm]m-4.south east);
    \node[mnota, anchor=west, align=left] at ([xshift=2.25cm]$(m-1.east)!0.5!(m-4.east)$)
      {16 byte consecutivi: 4 elementi\\da 4 byte ciascuno};
  \end{tikzpicture}
  \caption{\code{int a[4] = \{3, 7, 1, 9\};} in memoria (indirizzi di
           esempio).}
  \label{fig:array-memoria}
\end{figure}

\subsection{L'inizializzazione}

Gli array possono avere classe di memorizzazione automatica, esterna o
statica, ma non \code{register} (capitolo~\ref{chap:funzioni}). Nel C
tradizionale solo gli array esterni e statici potevano essere inizializzati
esplicitamente; in ANSI C lo possono essere anche quelli automatici, con
una lista di valori fra graffe:

\begin{itemize}
  \item \code{float v[5] = \{0.0, 1.0, 2.0, 3.0, 4.0\};} assegna
        \code{v[0] = 0.0}, \code{v[1] = 1.0}, e così via;
  \item \code{int v[5] = \{0\};} azzera l'intero vettore;
  \item \code{int v[5] = \{1\};} assegna \code{v[0] = 1} e pone a 0 i
        rimanenti: se i valori sono meno degli elementi, gli altri valgono 0;
  \item se la dimensione manca e c'è l'inizializzazione, la dimensione è il
        numero dei valori: \code{int a[] = \{1, 2, 3, 4, 5\};} equivale a
        \code{int a[5] = \{1, 2, 3, 4, 5\};}.
\end{itemize}

Se invece un array \textbf{non} viene inizializzato, dipende dalla sua
classe: gli array esterni e statici vengono azzerati dal sistema, quelli
automatici no, e contengono ciò che c'era in memoria.

\begin{affianco}[0.5][c]
\begin{lstlisting}[style=c, name=init.c]
#include <stdio.h>

int v[5];

int main(void)
{
   static int a[5];
   int b[5];
   int i;

   for (i=0; i<5; i++)
     printf("%d ",a[i]);
   printf("\n");

   for (i=0; i<5; i++)
     printf("%d ",v[i]);
   printf("\n");

   for (i=0; i<5; i++)
     printf("%d ",b[i]);
   printf("\n");
}
\end{lstlisting}
\risultato
\begin{console}
$ ./init
0 0 0 0 0
0 0 0 0 0
0 0 -2037142800 32652 179344864
$ ./init
0 0 0 0 0
0 0 0 0 0
0 0 744688368 32742 1091613072
\end{console}
\wtnota{\code{a} (statico) e \code{v} (esterno) sono azzerati; \code{b}
(automatico) contiene valori a caso, diversi a ogni esecuzione. \code{gcc}
con \code{-Wall} avvisa: \textit{‘b’ may be used uninitialized}.}
\end{affianco}

\subsection{Gli indici e l'errore più comune}

Se \code{a} è un array, \code{a[expr]}, con \code{expr} espressione intera,
è un suo elemento. Con la dichiarazione \code{int a[N];} gli indici validi
vanno da 0 a \code{N-1}. Accedere ad \code{a[i]} con \code{i} fuori da
questo intervallo è un errore di programmazione molto comune, e il C
\textbf{non lo controlla}: si legge o si scrive memoria che non è stata
riservata per l'array. L'effetto dipende dal sistema: il programma può
fermarsi con un errore, oppure proseguire con risultati sbagliati.

\begin{affianco}[0.5][c]
\begin{lstlisting}[style=c, name=fuori.c]
#include <stdio.h>

int main(void)
{
   int a[5] = {10, 20, 30, 40, 50};
   int i;

   for (i = 0; i <= 5; i++)      /* un giro di troppo */
      printf("a[%d] = %d\n", i, a[i]);
   return 0;
}
\end{lstlisting}
\risultato
\begin{console}
$ ./fuori
a[0] = 10
a[1] = 20
a[2] = 30
a[3] = 40
a[4] = 50
a[5] = 32764
\end{console}
\wtnota{\code{a[5]} non esiste: il programma legge senza protestare ciò che
c'è dopo l'array.}
\end{affianco}

Errori come questo sono difficili da trovare, perché spesso il programma
sembra funzionare. Un aiuto prezioso è l'opzione
\code{-fsanitize=address} di \code{gcc}, che aggiunge al programma dei
controlli sugli accessi alla memoria e lo ferma al primo accesso sbagliato,
dicendo dove (qui il rapporto è abbreviato, e al posto del percorso della
cartella c'è un segnaposto):

\begin{console}
$ gcc -g -fsanitize=address -o fuori fuori.c
$ ./fuori
a[0] = 10
a[1] = 20
a[2] = 30
a[3] = 40
a[4] = 50
=================================================================
==2607==ERROR: AddressSanitizer: stack-buffer-overflow on address 0x7f364ce00034 at pc 0x5575a6d3d467 bp 0x7fff693a7980 sp 0x7fff693a7970
READ of size 4 at 0x7f364ce00034 thread T0
    #0 0x5575a6d3d466 in main (*@\textit{[percorso del file]}@*)/fuori.c:9
(*@\textit{\dots}@*)
    [32, 52) 'a' (line 5) <== Memory access at offset 52 overflows this variable
(*@\textit{\dots}@*)
\end{console}

Il rapporto dice che alla riga 9 si è letto (\textit{READ}) un \code{int}
(4 byte) appena oltre la fine dell'array \code{a}, dichiarato alla riga 5.

\section{I puntatori}

\subsection{Indirizzi}

Una variabile è memorizzata in un certo numero di byte, in una particolare
posizione della memoria identificata da un \textbf{indirizzo}
(capitolo~\ref{chap:calcolatore}). L'indirizzo di una variabile \code{v} si
ottiene con l'operatore \code{\&}: \code{\&v}. Gli indirizzi sono valori, e
come tali si possono memorizzare e manipolare: le variabili che contengono
indirizzi si chiamano \term{puntatori}, e servono ad accedere alla memoria e
a manipolare gli indirizzi.

\subsection{Dichiarare e usare un puntatore}

La dichiarazione

\begin{lstlisting}[style=c]
int *p;
\end{lstlisting}

dichiara la variabile \code{p} come \textbf{puntatore a \code{int}}: può
contenere l'indirizzo di una variabile \code{int}. I suoi valori possibili
sono lo zero e gli interi positivi che sul sistema possono essere
interpretati come indirizzi. Alcuni esempi di assegnazione:

\begin{lstlisting}[style=c]
p = 0;
p = NULL;              /* equivalente a p = 0 */
p = &i;
p = (int *) 1776;      /* un indirizzo assoluto in memoria */
\end{lstlisting}

\code{NULL}, definita negli header standard, vale 0 e indica un puntatore
che \textbf{non punta a niente}. L'ultima assegnazione impone un indirizzo
preciso: non si fa mai nei programmi normali, perché quella cella
probabilmente non appartiene al programma.

L'operatore \code{*}, applicato a un puntatore, dà il \textbf{valore della
variabile puntata}: \code{*p} è il valore della variabile di cui \code{p}
contiene l'indirizzo (si dice che \code{p} viene \textit{dereferenziato}).
\code{*} e \code{\&} sono uno l'inverso dell'altro: \code{i} e \code{*\&i}
sono equivalenti. Seguiamo le istruzioni delle slide
(figura~\ref{fig:puntatore}):

\begin{lstlisting}[style=c]
int a, b, *p;

a = 1;
b = 2;
p = &a;
b = *p;               /* equivale a b = a */
\end{lstlisting}

\begin{figure}[H]
  \centering
  \begin{tikzpicture}[font=\Lato\scriptsize]
    \foreach \col/\testo in {0/{dopo \code{a = 1; b = 2;}}, 1/{dopo \code{p = \&a;}}, 2/{dopo \code{b = *p;}}}
      \node[font=\Lato\scriptsize\bfseries] at (\col*5.4+1.5,1.0) {\testo};
    \variabile{a1}{(1.4,0.3)}{a}{1}
    \variabile{b1}{(1.4,-0.3)}{b}{2}
    \variabile{p1}{(1.4,-0.9)}{p}{\mvuota}
    \begin{scope}[xshift=5.4cm]
      \variabile{a2}{(1.4,0.3)}{a}{1}
      \variabile{b2}{(1.4,-0.3)}{b}{2}
      \variabile{p2}{(1.4,-0.9)}{p}{}
      \riferimento{p2}{to[out=0, in=0, looseness=1.6] (a2.east)}
      \node[mnota, anchor=west] at (2.9,-0.9) {\code{p} contiene l'indirizzo di \code{a}};
    \end{scope}
    \begin{scope}[xshift=10.8cm]
      \variabile{a3}{(1.4,0.3)}{a}{1}
      \variabile[fill=accent!10!white]{b3}{(1.4,-0.3)}{b}{1}
      \variabile{p3}{(1.4,-0.9)}{p}{}
      \riferimento{p3}{to[out=0, in=0, looseness=1.6] (a3.east)}
      \node[mnota, anchor=west, align=left] at (2.9,-0.3) {\code{*p} è il valore\\di \code{a}: 1};
    \end{scope}
  \end{tikzpicture}
  \caption{Un puntatore si disegna come una freccia verso la variabile di
           cui contiene l'indirizzo.}
  \label{fig:puntatore}
\end{figure}

\subsection{Che cosa contiene un puntatore}

Un indirizzo si stampa con la conversione \code{\%p}, in esadecimale:

\begin{affianco}[0.5][c]
\begin{lstlisting}[style=c, name=indirizzo.c]
#include <stdio.h>

int main(void)
{
   int   i = 7, *p;

   p = &i;
   printf("%s%d\n%s%p\n",
      "   Value of i: ", *p,
      "Location of i: ",  p);
   return 0;
}
\end{lstlisting}
\risultato
\begin{console}
$ ./indirizzo
   Value of i: 7
Location of i: 0x7ffda49df4ac
$ ./indirizzo
   Value of i: 7
Location of i: 0x7ffe989ec3bc
\end{console}
\end{affianco}

L'indirizzo cambia a ogni esecuzione: per sicurezza, il sistema operativo
colloca la memoria del programma in posizioni diverse ogni volta. Nelle
slide le due stringhe contengono anche un \code{\%d}
(\code{"   Value of i: \%d"}): così, però, \code{\%d} non è una conversione
(sta in un argomento, non nella stringa di controllo) e viene stampato tale
e quale, \code{   Value of i: \%d7}. La versione qui sopra è quella corretta.

\subsection{Espressioni con i puntatori}

Con le dichiarazioni

\begin{lstlisting}[style=c]
int i = 3, j = 5, *p = &i, *q = &j, *r;
double x;
\end{lstlisting}

ecco alcune espressioni, con l'espressione equivalente e il valore
(verificato con il compilatore, calcolando ciascuna partendo da questi
valori):

\begin{center}
\renewcommand{\arraystretch}{1.15}
\begin{tabular}{@{}>{\FiraCodeNL}l>{\FiraCodeNL}l l@{}}
  \toprule
  \normalfont\textbf{espressione} & \normalfont\textbf{equivalente} & \textbf{valore} \\
  \midrule
  p == \&i            & p == (\&i)                     & 1 \\
  **\&p               & *(*(\&p))                      & 3 \\
  *(r = \&j) *= *p    & (*(r = (\&j))) *= *p           & 15 \\
  7 * *p / *q + 7     & (((7 * (*p))) / (*q)) + 7      & 11 \\
  r = \&x             & r = \&x                        & illegale \\
  \bottomrule
\end{tabular}
\end{center}

\begin{itemize}
  \item \code{**\&p}: \code{\&p} è l'indirizzo di \code{p}; \code{*} lo
        riporta a \code{p}, e un altro \code{*} dà il valore puntato da
        \code{p}, cioè \code{i}: 3.
  \item \code{*(r = \&j) *= *p}: \code{r} riceve l'indirizzo di \code{j}, e
        \code{*r *= *p} moltiplica \code{j} per \code{i}: \code{j} diventa 15.
  \item \code{7 * *p / *q + 7}: $(7 \cdot 3) / 5 + 7 = 21 / 5 + 7 = 4 + 7 =
        11$ (divisione intera). Nelle slide è scritta \code{7**p / *q+7}: gli
        spazi non contano, ma rendono leggibile.
  \item \code{r = \&x} è illegale: \code{r} punta a \code{int}, \code{\&x} è
        l'indirizzo di un \code{double}. \code{gcc} lo segnala (e i
        compilatori più recenti lo considerano un errore):
\end{itemize}

\begin{console}
$ gcc -Wall -o rx rx.c
rx.c: In function (*@‘@*)main(*@’@*):
rx.c:9:5: warning: assignment to (*@‘@*)int *(*@’@*) from incompatible pointer type (*@‘@*)double *(*@’@*) [-Wincompatible-pointer-types]
    9 |   r = &x;
      |     ^
\end{console}

\subsection{Che cosa non si può puntare}

Non tutto ha un indirizzo. Non si può puntare a una \textbf{costante}
(\code{\&3}), a un'\textbf{espressione} (\code{\&(k + 9)}) né a una
variabile \code{register} (\code{\&v}). Le prime due non sono celle di
memoria con un nome (sono \textit{lvalue} mancati); la terza, come abbiamo
visto, sta in un registro:

\begin{console}
$ gcc -Wall -o nonpuntabili nonpuntabili.c
nonpuntabili.c: In function (*@‘@*)main(*@’@*):
nonpuntabili.c:6:8: error: lvalue required as unary (*@‘@*)&(*@’@*) operand
nonpuntabili.c:7:8: error: lvalue required as unary (*@‘@*)&(*@’@*) operand
nonpuntabili.c:8:4: error: address of register variable (*@‘@*)v(*@’@*) requested
\end{console}

Lo stesso vale per un cast: nella cartella degli esempi c'è un tentativo,
\code{r = \&((int) x);}, di far puntare un \code{int *} a un
\code{double} “convertito”. Il risultato di un cast è un valore, non una
variabile, e il compilatore risponde ancora \textit{lvalue required as
unary ‘\&’ operand}.

\section{Il passaggio per indirizzo}

\subsection{Modificare le variabili del chiamante}

Una variabile passata come parametro a una funzione non cambia valore dopo
la chiamata: è il passaggio \textbf{per valore}
(capitolo~\ref{chap:funzioni}). Altri linguaggi hanno anche il passaggio
\textbf{per riferimento}, in cui la funzione può modificare la variabile
passata. In C si ottiene un effetto simile con i puntatori: si passa
l'\textbf{indirizzo} della variabile, e la funzione lavora sulla variabile
attraverso \code{*}.

\begin{affianco}[0.5][t]
\begin{lstlisting}[style=c, name=valore.c]
int f(int a, int b)
{
  a = a+1;
  b = b+1;
  return (a*b);
}

int main()
{
  int a = 1, b = 2, c;

  c = f(a,b);
  printf("a = %d\nb = %d\nc = %d \n", a, b, c);
}
\end{lstlisting}
\begin{console}
$ ./valore
a = 1
b = 2
c = 6
\end{console}
\risultato
\begin{lstlisting}[style=c, name=indirizzi.c]
int f(int *a, int *b)
{
  *a = *a + 1;
  *b = *b + 1;
  return (*a * *b);
}

int main()
{
  int a = 1, b = 2, c;

  c = f(&a, &b);
  printf("a = %d\nb = %d\nc = %d \n", a, b, c);
}
\end{lstlisting}
\begin{console}
$ ./indirizzi
a = 2
b = 3
c = 6
\end{console}
\end{affianco}

A sinistra \code{f} incrementa le proprie copie; a destra riceve gli
indirizzi delle variabili di \code{main} e incrementa \textbf{quelle}
(figura~\ref{fig:per-indirizzo}). Si noti che \code{a} e \code{b}, dentro
\code{f}, ora sono puntatori: \code{*a} è l'intero puntato.

\begin{figure}[H]
  \centering
  \begin{tikzpicture}[font=\Lato\scriptsize]
    \variabile{ma}{(1.6,0.0)}{a}{1 \textcolor{accent!75!black}{$\rightarrow$ 2}}
    \variabile{mb}{(1.6,-0.6)}{b}{2 \textcolor{accent!75!black}{$\rightarrow$ 3}}
    \variabile{mc}{(1.6,-1.2)}{c}{6}
    \scopo{2}{m}{main}{(ma-nome) (ma) (mb-nome) (mb) (mc-nome) (mc)}
    \node[mvar, fill=accent!10!white] (fa) at (6.4,-0.2) {};
    \node[mnome, anchor=west] (fa-nome) at ([xshift=2pt]fa.east) {a};
    \node[mvar, fill=accent!10!white] (fb) at (6.4,-0.8) {};
    \node[mnome, anchor=west] (fb-nome) at ([xshift=2pt]fb.east) {b};
    \scopo{3}{f}{f}{(fa-nome) (fa) (fb-nome) (fb)}
    \riferimento{fa}{-- (ma.east)}
    \riferimento{fb}{-- (mb.east)}
    \node[mnota, anchor=west, align=left] at (7.9,-0.5) {i parametri di \code{f} contengono gli indirizzi:\\\code{*a = *a + 1} incrementa la \code{a} di \code{main}};
  \end{tikzpicture}
  \caption{Il passaggio per indirizzo: \code{f} arriva alle variabili di
           \code{main} attraverso i puntatori.}
  \label{fig:per-indirizzo}
\end{figure}

\warning{Una svista nelle slide}{
  Nella versione delle slide il corpo di \code{f} è rimasto quello di
  prima: \code{a = a+1; b = b+1; return (a*b);}. Ma ora \code{a} e
  \code{b} sono puntatori: \code{a = a+1} sposterebbe il puntatore, non
  incrementerebbe il valore, e moltiplicare due puntatori non ha senso. Il
  compilatore infatti si ferma: \textit{invalid operands to binary *
  (have ‘int *’ and ‘int *’)}. Bisogna usare \code{*} ovunque si intenda il
  valore puntato, come nella versione qui sopra.}

\subsection{Lo scambio di due variabili}

Lo scambio (\textit{swap}) dei valori di due variabili si fa di solito con
una variabile di appoggio: \code{tmp = a; a = b; b = tmp;}. Una funzione che
scambi due variabili del chiamante deve riceverne gli indirizzi:

\begin{affianco}[0.5][c]
\begin{lstlisting}[style=c, name=swap.c]
#include <stdio.h>

void   swap(int *, int *);

int main(void)
{
   int    i = 3, j = 5;

   swap(&i, &j);
   printf("%d  %d\n", i, j);       /* 5  3 is printed */
   return 0;
}

void swap(int *p, int *q)
{
   int   tmp;

   tmp = *p;
   *p = *q;
   *q = tmp;
}
\end{lstlisting}
\risultato
\begin{console}
$ ./swap
5  3
\end{console}
\end{affianco}

\begin{figure}[H]
  \centering
  \begin{tikzpicture}[font=\Lato\scriptsize]
    \foreach \col/\testo/\vi/\vj/\vt in {0/{\code{tmp = *p;}}/3/5/3, 1/{\code{*p = *q;}}/5/5/3, 2/{\code{*q = tmp;}}/5/3/3} {
      \begin{scope}[xshift=\col*5.6cm]
        \node[font=\Lato\scriptsize\bfseries] at (2.2,1.0) {dopo \testo};
        \variabile{i\col}{(1.0,0.3)}{i}{\vi}
        \variabile{j\col}{(1.0,-0.3)}{j}{\vj}
        \node[mvar, fill=accent!10!white] (p\col) at (3.6,0.3) {};
        \node[mnome, anchor=west] at ([xshift=2pt]p\col.east) {p};
        \node[mvar, fill=accent!10!white] (q\col) at (3.6,-0.3) {};
        \node[mnome, anchor=west] at ([xshift=2pt]q\col.east) {q};
        \node[mvar, fill=accent!10!white] (t\col) at (3.6,-0.9) {\vt};
        \node[mnome, anchor=west] at ([xshift=2pt]t\col.east) {tmp};
        \riferimento{p\col}{-- (i\col.east)}
        \riferimento{q\col}{-- (j\col.east)}
      \end{scope}
    }
  \end{tikzpicture}
  \caption{\code{swap} passo per passo: \code{p} e \code{q} puntano a
           \code{i} e \code{j} di \code{main}.}
  \label{fig:swap}
\end{figure}

È anche il motivo per cui \code{scanf} riceve \code{\&x}: per scrivere in
\code{x}, deve sapere dove si trova (capitolo~\ref{chap:primi-programmi}).

\section{Array e puntatori}

\subsection{Il nome di un array è un indirizzo}

Il nome di un array è un indirizzo, cioè il valore di un puntatore: quello
del suo primo elemento. Da qui le relazioni fondamentali:

\begin{itemize}
  \item se \code{a} è un array e \code{i} un intero, \code{a[i]} (l'elemento
        di indice \code{i}) è equivalente a \code{*(a + i)} (il contenuto
        dell'indirizzo \code{a + i});
  \item se \code{p} è un puntatore, anche \code{p[i]} è equivalente a
        \code{*(p + i)}: i puntatori si possono indicizzare come gli array;
  \item con \code{int a[N], *p;}, l'istruzione \code{p = a;} equivale a
        \code{p = \&a[0];}, e \code{p = a + 1;} equivale a
        \code{p = \&a[1];}.
\end{itemize}

Il punto chiave è l'\textbf{aritmetica dei puntatori}: sommare 1 a un
puntatore non lo sposta di un byte, ma di \textbf{un elemento}, cioè di
tanti byte quanti ne occupa il tipo puntato (4 per un \code{int}).

\begin{affianco}[0.55][c]
\begin{lstlisting}[style=c, name=aritmetica.c]
#include <stdio.h>

int main(void)
{
   int a[3] = {10, 20, 30};
   int *p = a;

   printf("a     = %p\n", (void *) a);
   printf("&a[0] = %p\n", (void *) &a[0]);
   printf("a + 1 = %p\n", (void *) (a + 1));
   printf("&a[1] = %p\n", (void *) &a[1]);
   printf("a + 2 = %p\n", (void *) (a + 2));
   printf("*(a + 2) = %d   p[2] = %d   *(p + 2) = %d\n",
          *(a + 2), p[2], *(p + 2));
   return 0;
}
\end{lstlisting}
\risultato
\begin{console}
$ ./aritmetica
a     = 0x7ffc5944463c
&a[0] = 0x7ffc5944463c
a + 1 = 0x7ffc59444640
&a[1] = 0x7ffc59444640
a + 2 = 0x7ffc59444644
*(a + 2) = 30   p[2] = 30   *(p + 2) = 30
\end{console}
\wtnota{Da \code{a} ad \code{a + 1} l'indirizzo cresce di 4
(\code{...63c}, \code{...640}): un \code{int}.}
\end{affianco}

\begin{figure}[H]
  \centering
  \begin{tikzpicture}[font=\FiraCodeNL\scriptsize]
    \noindici
    \celle[minimum width=2.0cm]{c}{(0,0)}{10,20,30}
    \foreach \k/\ind in {0/1000, 1/1004, 2/1008}
      \node[mindir, anchor=north] at ([yshift=-1pt]c-\k.south) {\ind};
    \foreach \k/\testo in {0/{a\\\&a[0]\\p}, 1/{a + 1\\\&a[1]}, 2/{a + 2\\\&a[2]}} {
      \node[font=\FiraCodeNL\scriptsize, text=mref, anchor=south, align=center] (l\k) at ([yshift=0.5cm]c-\k.north) {\testo};
      \draw[mrif] (l\k.south) -- (c-\k.north);
    }
    \node[mnota, anchor=west, align=left] at ([xshift=0.5cm]c-2.east) {\code{a[i]} è \code{*(a + i)}:\\il contenuto della cella \code{a + i}};
  \end{tikzpicture}
  \caption{\code{a + i} è l'indirizzo dell'elemento \code{i} (indirizzi di
           esempio): ogni passo vale 4 byte.}
  \label{fig:aritmetica}
\end{figure}

\subsection{Quattro modi di sommare un array}

Grazie a queste equivalenze, la somma degli elementi di un array si può
scrivere in molti modi, tutti equivalenti:

\begin{affianco}[0.55][c]
\begin{lstlisting}[style=c, name=quattro.c]
#include <stdio.h>

#define N 5

int main(void)
{
   int a[N] = {1, 2, 3, 4, 5}, i, *p, sum;

   sum = 0;
   for (i = 0; i < N; i++)
      sum += a[i];
   printf("a[i]:      %d\n", sum);

   sum = 0;
   for (p = a, i = 0; i < N; i++)
      sum += p[i];
   printf("p[i]:      %d\n", sum);

   sum = 0;
   for (i = 0; i < N; i++)
      sum += *(a + i);
   printf("*(a + i):  %d\n", sum);

   sum = 0;
   for (p = a; p < &a[N]; p++)
      sum += *p;
   printf("*p, p++:   %d\n", sum);
   return 0;
}
\end{lstlisting}
\risultato
\begin{console}
$ ./quattro
a[i]:      15
p[i]:      15
*(a + i):  15
*p, p++:   15
\end{console}
\vspace{2pt}
\wtnota{Nell'ultimo ciclo è il puntatore a scorrere l'array: parte da
\code{a} e avanza di un elemento a ogni giro, finché arriva a
\code{\&a[N]}, l'indirizzo subito dopo l'ultimo elemento.}
\end{affianco}

\subsection{Un array non si può spostare}

C'è però una differenza fra un array e un puntatore: un puntatore può
assumere indirizzi diversi, mentre il nome di un array è un indirizzo
\textbf{costante}, che non si può cambiare. Le istruzioni \code{a = p;},
\code{++a;} e \code{a += 2;} sono illegali:

\begin{console}
$ gcc -Wall -o costante costante.c
costante.c: In function (*@‘@*)main(*@’@*):
costante.c:5:6: error: assignment to expression with array type
costante.c:6:4: error: lvalue required as increment operand
costante.c:7:6: error: assignment to expression with array type
\end{console}

\section{Gli array come parametri}

\subsection{Si passa l'indirizzo}

Nella definizione di una funzione, un parametro dichiarato come array è in
realtà un \textbf{puntatore}. Quando si passa un array come argomento, viene
passato per valore il suo \textbf{indirizzo base}: gli elementi
\textbf{non vengono copiati}. Come esempio, una funzione che somma gli
elementi di un array di \code{double}: le due intestazioni seguenti sono
equivalenti.

\begin{affianco}[0.5][c]
\begin{lstlisting}[style=c]
double sum_array(double a[], int n)
\end{lstlisting}
\risultato
\begin{lstlisting}[style=c]
double sum_array(double *a, int n)
\end{lstlisting}
\end{affianco}

Siccome la funzione riceve solo un indirizzo, non sa quanto è lungo
l'array: la dimensione va passata come parametro (qui \code{n}). E siccome
riceve un indirizzo qualsiasi, può lavorare anche su una \textbf{parte}
dell'array. Con \code{double v[100];} e \code{v[i]} uguale a \code{i}:

\begin{affianco}[0.6][c]
\codicepiccolo
\begin{lstlisting}[style=c, name=sum_array.c]
#include <stdio.h>

double sum_array(double a[], int n)   /* n e' la dimensione di a[] */
{
   int    i;
   double sum = 0.0;

   for (i = 0; i < n; ++i)
      sum += a[i];
   return sum;
}

int main(void)
{
   double v[100];
   int    i, k = 10;

   for (i = 0; i < 100; i++)
      v[i] = i;                       /* v[i] vale i */
   printf("%.0f\n", sum_array(v, 100));        /* v[0] + ... + v[99]   */
   printf("%.0f\n", sum_array(v, 88));         /* v[0] + ... + v[87]   */
   printf("%.0f\n", sum_array(&v[7], k - 7));  /* v[7] + ... + v[k-1]  */
   printf("%.0f\n", sum_array(v + 7, 2 * k));  /* v[7] + ... + v[2k+6] */
   return 0;
}
\end{lstlisting}
\risultato
\begin{console}
$ ./sum_array
4950
3828
24
330
\end{console}
\wtnota{Con \code{k} uguale a 10: $7 + 8 + 9 = 24$, e la somma da 7 a 26
vale 330.}
\end{affianco}

Nella chiamata \code{sum\_array(v + 7, 2 * k)} il parametro \code{a} punta
a \code{v[7]}: per la funzione \code{a[0]} è \code{v[7]}, \code{a[1]} è
\code{v[8]}, e così via (figura~\ref{fig:parte-array}).

\begin{figure}[H]
  \centering
  \begin{tikzpicture}[font=\FiraCodeNL\scriptsize]
    \noindici
    \celle[minimum width=0.9cm]{v}{(0,0)}{0,\ldots,6,7,8,9,\ldots,26,27,\ldots}
    \foreach \k/\ind in {0/v[0], 2/v[6], 3/v[7], 4/v[8], 5/v[9], 7/v[26], 8/v[27]}
      \node[mindir, anchor=north] at ([yshift=-1pt]v-\k.south) {\ind};
    \foreach \k/\ind in {3/a[0], 4/a[1], 5/a[2], 7/a[19]}
      \node[font=\FiraCodeNL\tiny, text=accent!75!black, anchor=south] at ([yshift=1pt]v-\k.north) {\ind};
    \draw[wtbrace, draw=accent!70!black] ([yshift=-0.35cm]v-3.south west) -- ([yshift=-0.35cm]v-7.south east);
    \node[mnota, anchor=north, text=accent!70!black] at ([yshift=-0.65cm]$(v-3.south)!0.5!(v-7.south)$) {i 20 elementi sommati da \code{sum\_array(v + 7, 20)}};
    \node[font=\FiraCodeNL\scriptsize, text=mref, anchor=south] (pa) at ([yshift=0.8cm]v-3.north) {a};
    \draw[mrif] (pa.south) -- ([yshift=0.3cm]v-3.north);
  \end{tikzpicture}
  \caption{La funzione vede l'array a partire dall'indirizzo che riceve.}
  \label{fig:parte-array}
\end{figure}

\subsection{\code{sizeof} su un parametro}

Una conseguenza che sorprende: dentro la funzione, \code{sizeof} applicato
al parametro dà la dimensione di un \textbf{puntatore}, non dell'array.

\begin{affianco}[0.55][c]
\begin{lstlisting}[style=c, name=dimensione.c]
#include <stdio.h>

void stampa_dimensione(int a[])
{
   printf("nella funzione: sizeof(a) = %zu\n", sizeof(a));
}

int main(void)
{
   int v[10];

   printf("in main:        sizeof(v) = %zu\n", sizeof(v));
   stampa_dimensione(v);
   return 0;
}
\end{lstlisting}
\risultato
\begin{console}
$ ./dimensione
in main:        sizeof(v) = 40
nella funzione: sizeof(a) = 8
\end{console}
\wtnota{40 byte sono i dieci \code{int}; 8 byte sono un indirizzo.
\code{gcc} avvisa: \textit{‘sizeof’ on array function parameter ‘a’ will
return size of ‘int *’}.}
\end{affianco}

\section{Ordinare un array}

Un'applicazione classica degli array (e dei puntatori, per gli scambi) è
l'\textbf{ordinamento}: disporre gli elementi in ordine crescente. Le
slide presentano due algoritmi semplici. Entrambi usano la \code{swap} vista
sopra, a cui passano gli indirizzi di due elementi, \code{a + j - 1} e
\code{a + j}, cioè \code{\&a[j-1]} e \code{\&a[j]}.

\subsection{Il bubble sort}

\begin{affianco}[0.55][c]
\begin{lstlisting}[style=c, name=bubblesort.c]
void bubble(int *a, int n)     /* n is the size of a */
{
   int    i, j;

   for (i = 0; i < n - 1; ++i)
      for (j = n - 1; j > i; --j)
         if (a[j-1] > a[j])
            swap(a+j-1, a+j);
}
\end{lstlisting}
\risultato
\begin{console}
$ ./bubblesort
7 3 66 3 -5 22 -77 2
-77 -5 2 3 3 7 22 66
\end{console}
\wtnota{\code{main} stampa l'array prima e dopo la chiamata
\code{bubble(v, 8)}.}
\end{affianco}

A ogni passo (un giro del ciclo esterno), il ciclo interno scorre l'array
\textbf{dal fondo verso l'inizio} confrontando ogni coppia di elementi
vicini, e li scambia se sono in ordine sbagliato. Così il più piccolo degli
elementi rimasti “risale” come una bolla fino alla posizione \code{i}:
dopo il passo $k$, i primi $k$ elementi sono definitivamente al loro posto
(figura~\ref{fig:bubble}).

\begin{figure}[H]
  \centering
  \setlength{\tabcolsep}{6pt}
  \renewcommand{\arraystretch}{1.12}
  \newcommand{\ok}[1]{\cellcolor{green-gradient-1!45!white}#1}
  \begin{tabular}{@{}l>{\FiraCodeNL}r>{\FiraCodeNL}r>{\FiraCodeNL}r>{\FiraCodeNL}r>{\FiraCodeNL}r>{\FiraCodeNL}r>{\FiraCodeNL}r>{\FiraCodeNL}r@{}}
    \toprule
    Dati    & 7 & 3 & 66 & 3 & -5 & 22 & -77 & 2 \\
    Passo 1 & \ok{-77} & 7 & 3 & 66 & 3 & -5 & 22 & 2 \\
    Passo 2 & \ok{-77} & \ok{-5} & 7 & 3 & 66 & 3 & 2 & 22 \\
    Passo 3 & \ok{-77} & \ok{-5} & \ok{2} & 7 & 3 & 66 & 3 & 22 \\
    Passo 4 & \ok{-77} & \ok{-5} & \ok{2} & \ok{3} & 7 & 3 & 66 & 22 \\
    Passo 5 & \ok{-77} & \ok{-5} & \ok{2} & \ok{3} & \ok{3} & 7 & 22 & 66 \\
    Passo 6 & \ok{-77} & \ok{-5} & \ok{2} & \ok{3} & \ok{3} & \ok{7} & 22 & 66 \\
    Passo 7 & \ok{-77} & \ok{-5} & \ok{2} & \ok{3} & \ok{3} & \ok{7} & \ok{22} & \ok{66} \\
    \bottomrule
  \end{tabular}
  \caption{Il bubble sort passo per passo (verificato stampando l'array dopo
           ogni giro del ciclo esterno): in verde la parte già ordinata.}
  \label{fig:bubble}
\end{figure}

Gli ultimi passi non cambiano più nulla: l'array è già ordinato dal passo 5,
ma l'algoritmo non se ne accorge e fa comunque tutti gli $n - 1$ passi.

\subsection{Il selection sort}

\begin{affianco}[0.55][c]
\begin{lstlisting}[style=c, name=selectionsort.c]
void selection_sort(int *a, int n)     /* n is the size of a */
{
   int    i, j;

   for (i = 0; i < n - 1; ++i)
      for (j = i + 1; j < n ; j++)
         if (a[j] < a[i])
            swap(a+j, a+i);
}
\end{lstlisting}
\risultato
\begin{console}
$ ./selectionsort
7 3 66 3 -5 22 -77 2
-77 -5 2 3 3 7 22 66
\end{console}
\end{affianco}

Qui al passo \code{i} si confronta l'elemento in posizione \code{i} con
tutti quelli che seguono, scambiandolo ogni volta che se ne trova uno più
piccolo: alla fine del passo, in posizione \code{i} c'è il più piccolo degli
elementi rimasti, cioè quello che va lì (figura~\ref{fig:selection}).

\begin{figure}[H]
  \centering
  \setlength{\tabcolsep}{6pt}
  \renewcommand{\arraystretch}{1.12}
  \newcommand{\ok}[1]{\cellcolor{green-gradient-1!45!white}#1}
  \begin{tabular}{@{}l>{\FiraCodeNL}r>{\FiraCodeNL}r>{\FiraCodeNL}r>{\FiraCodeNL}r>{\FiraCodeNL}r>{\FiraCodeNL}r>{\FiraCodeNL}r>{\FiraCodeNL}r@{}}
    \toprule
    Dati    & 7 & 3 & 66 & 3 & -5 & 22 & -77 & 2 \\
    Passo 1 & \ok{-77} & 7 & 66 & 3 & 3 & 22 & -5 & 2 \\
    Passo 2 & \ok{-77} & \ok{-5} & 66 & 7 & 3 & 22 & 3 & 2 \\
    Passo 3 & \ok{-77} & \ok{-5} & \ok{2} & 66 & 7 & 22 & 3 & 3 \\
    Passo 4 & \ok{-77} & \ok{-5} & \ok{2} & \ok{3} & 66 & 22 & 7 & 3 \\
    Passo 5 & \ok{-77} & \ok{-5} & \ok{2} & \ok{3} & \ok{3} & 66 & 22 & 7 \\
    Passo 6 & \ok{-77} & \ok{-5} & \ok{2} & \ok{3} & \ok{3} & \ok{7} & 66 & 22 \\
    Passo 7 & \ok{-77} & \ok{-5} & \ok{2} & \ok{3} & \ok{3} & \ok{7} & \ok{22} & \ok{66} \\
    \bottomrule
  \end{tabular}
  \caption{Il selection sort passo per passo: a ogni passo il minimo dei
           rimanenti va nella sua posizione definitiva.}
  \label{fig:selection}
\end{figure}

\section{L'allocazione dinamica}

\subsection{\code{malloc}, \code{calloc} e \code{free}}

Finora la dimensione di un array era una costante, decisa scrivendo il
programma. Spesso però la si conosce solo durante l'esecuzione (quanti
numeri inserirà l'utente?). Per questo il C fornisce le funzioni
\code{malloc} (\textit{memory allocation}) e \code{calloc}
(\textit{contiguous allocation}), con i prototipi in \code{stdlib.h}: creano
spazio \textbf{dinamicamente}, mentre il programma è in esecuzione, per
array, strutture e unioni.

\begin{itemize}
  \item \code{calloc(n, el\_size)} riserva spazio contiguo per un array di
        \code{n} elementi da \code{el\_size} byte ciascuno (la dimensione
        si ottiene con \code{sizeof}), e lo inizializza con tutti i bit a
        zero. I due parametri sono di tipo \code{size\_t}.
  \item \code{malloc(size)} riserva \code{size} byte e non li inizializza:
        può essere un po' più veloce.
  \item Se riescono, entrambe restituiscono un puntatore di tipo
        \code{void *} (un indirizzo “generico”) alla base della memoria;
        altrimenti restituiscono \code{NULL}.
\end{itemize}

\begin{lstlisting}[style=c]
a = (int *) calloc(n, sizeof(int));
a = (int *) malloc(n * sizeof(int));
\end{lstlisting}

Il cast \code{(int *)} si può anche omettere, perché in C un
\code{void *} si converte da solo in qualsiasi altro puntatore. Lo spazio
allocato dinamicamente \textbf{non viene liberato} automaticamente
all'uscita dalle funzioni: va restituito esplicitamente con \code{free(ptr)}.
Se \code{ptr} è \code{NULL}, \code{free} non fa nulla; altrimenti
\code{ptr} deve essere esattamente l'indirizzo base restituito da
\code{malloc} o \code{calloc}, o il comportamento è indefinito (e l'effetto
dipende dal sistema).

\subsection{Un esempio}

Il programma seguente chiede ripetutamente un numero \code{n}, alloca un
array di \code{n} interi, lo riempie di numeri casuali fra $-9$ e 9, lo
stampa con la somma e lo libera. \code{rand()} (in \code{stdlib.h})
restituisce un intero pseudo-casuale, e \code{srand(time(NULL))} sceglie il
seme in base all'ora, così che i numeri cambino a ogni esecuzione.

\begin{lstlisting}[style=c, name=dinamico.c]
#include <stdio.h>
#include <stdlib.h>
#include <time.h>

void   fill_array(int *a, int n);
int    sum_array(int *a, int n);
void   wrt_array(int *a, int n);

int main(void)
{
   int   *a, n;

   srand(time(NULL));     /* seed the random number generator */
   printf("Input n: ");
   while (scanf("%d", &n) == 1 && n >= 1) {
      putchar('\n');
      a = calloc(n, sizeof(int));     /* allocate space for a[] */
      fill_array(a, n);
      wrt_array(a, n);
      printf("sum = %d\n\n", sum_array(a, n));
      free(a);
      printf("Input n: ");
   }
   printf("\nBye!\n\n");
   return 0;
}

void fill_array(int *a, int n)
{
   int   i;

   for (i = 0; i < n; ++i)
      a[i] = rand() % 19 - 9;
}

int sum_array(int *a, int n)
{
   int   i, sum = 0;

   for (i = 0; i < n; ++i)
      sum += a[i];
   return sum;
}

void wrt_array(int *a, int n)
{
   int   i;

   printf("a = [");
   for (i = 0; i < n; ++i)
      printf("%d%s", a[i], ((i < n - 1) ? ", " : "]\n"));
}
\end{lstlisting}

\subsection{Il programma in esecuzione}

Ecco un'esecuzione (i numeri, naturalmente, cambiano ogni volta):

\begin{affianco}[0.42][c]
\begin{console}
$ ./dinamico
Input n: <<5>>

a = [8, -1, -5, -9, -3]
sum = -10

Input n: <<3>>

a = [-4, -1, 7]
sum = 2

Input n: <<0>>

Bye!
\end{console}
\risultato
\begin{itemize}
  \item \code{rand() \% 19} vale da 0 a 18; togliendo 9 si va da $-9$ a 9.
  \item In \code{wrt\_array}, \code{(i < n - 1) ? ", " : "]\textbackslash
        n"} è l'\textbf{operatore condizionale}: vale \code{", "} se la
        condizione è vera, \code{"]\textbackslash n"} altrimenti. Così dopo
        l'ultimo elemento si chiude la parentesi invece di mettere la
        virgola.
  \item Il ciclo continua finché \code{scanf} legge un numero
        \textbf{e} il numero è almeno 1.
\end{itemize}
\end{affianco}

\warning{Una svista nelle slide}{
  Nelle slide la condizione è \code{scanf("\%d", \&n) == 1 || n >= 1}, con
  l'\textit{or}. Così, se si inserisce 0, la prima parte è vera e il ciclo
  prosegue con un array vuoto; e se l'input finisce (\tasti{Ctrl,D}) dopo un
  numero valido, la prima parte è falsa ma \code{n} è rimasto quello di
  prima, la seconda è vera, e il ciclo non termina più (in una prova ha
  stampato oltre 10 milioni di righe in un secondo). La condizione giusta,
  quella del libro da cui viene l'esempio, è con \code{\&\&}.}

\subsection{Lo stack e lo heap}

Dove finisce la memoria allocata dinamicamente? La memoria di un programma
in esecuzione è divisa in zone (figura~\ref{fig:zone-memoria}):

\begin{itemize}
  \item lo \term{stack} (la pila) è lo spazio di lavoro delle funzioni:
        a ogni chiamata viene riservato un blocco in cima allo stack per le
        variabili locali e alcuni dati, e viene tolto al ritorno
        (figura~\ref{fig:pila});
  \item lo \term{heap} è la memoria riservata all'allocazione dinamica:
        \code{malloc} e \code{calloc} prendono spazio da lì, e \code{free}
        lo restituisce;
  \item le variabili globali e \code{static} hanno una loro zona, che esiste
        per tutta l'esecuzione, e il codice del programma un'altra ancora.
\end{itemize}

Lo stack è tipicamente molto più piccolo dello heap (su Linux, di solito, 8
MB): un array enorme dichiarato come variabile locale può esaurirlo
(\textit{stack overflow}), mentre lo stesso array allocato con
\code{malloc} non dà problemi.

\begin{figure}[H]
  \centering
  \begin{tikzpicture}[font=\Lato\scriptsize,
      zona/.style={draw=primary!60!black, line width=0.7pt, minimum width=4.2cm, align=center}]
    \node[zona, fill=accent!10!white, minimum height=1.3cm] (st) at (0,2.7) {\textbf{stack}\\variabili locali, parametri};
    \node[zona, fill=white, minimum height=1.2cm, draw=neutral!40!white, dashed] (lib) at (0,1.45) {\color{neutral!45!black}spazio libero};
    \node[zona, fill=green-gradient-1!40!white, minimum height=1.3cm] (hp) at (0,0.2) {\textbf{heap}\\\code{malloc}, \code{calloc}};
    \node[zona, fill=primary!8!white, minimum height=0.75cm] (gl) at (0,-0.83) {variabili globali e \code{static}};
    \node[zona, fill=neutral!10!white, minimum height=0.75cm] (cd) at (0,-1.6) {codice del programma};
    \draw[wtarrowacc] ([xshift=-1.5cm]st.south) -- ++(0,-0.45) node[left, mnota, pos=0.3] {cresce verso il basso};
    \draw[wtarrowacc] ([xshift=-1.5cm]hp.north) -- ++(0,0.45) node[left, mnota, pos=0.3] {cresce verso l'alto};
    \node[mnota, anchor=west] at (2.4,2.7) {un blocco per ogni chiamata, tolto al ritorno};
    \node[mnota, anchor=west] at (2.4,0.2) {resta finché non si chiama \code{free}};
    \node[mnota, anchor=west] at (2.4,-1.2) {esistono per tutta l'esecuzione};
  \end{tikzpicture}
  \caption{Le zone della memoria di un programma (lo schema è semplificato).}
  \label{fig:zone-memoria}
\end{figure}

\insintesi{
  Un \textbf{array} raccoglie elementi dello stesso tipo in celle
  consecutive; gli indici vanno da 0 a $N-1$, e il C non controlla gli
  accessi fuori dai limiti (\code{-fsanitize=address} aiuta a trovarli). Gli
  array automatici non inizializzati contengono valori casuali. Un
  \textbf{puntatore} contiene un indirizzo: \code{\&v} è l'indirizzo di
  \code{v}, \code{*p} il valore puntato, \code{NULL} il puntatore che non
  punta a niente. Passando gli indirizzi una funzione può modificare le
  variabili del chiamante (\code{swap}). Il nome di un array è l'indirizzo
  costante del primo elemento: \code{a[i]} equivale a \code{*(a + i)}, e
  \code{p + 1} avanza di un elemento. Un array passato a una funzione è
  passato come indirizzo (la dimensione va passata a parte). Bubble sort e
  selection sort ordinano con confronti e scambi. \code{malloc} e
  \code{calloc} allocano memoria nello heap durante l'esecuzione, e
  \code{free} la restituisce.}

\section{Riferimenti}

\begin{itemize}
  \item \textbf{A. Kelley, I. Pohl}, \textit{A Book on C}: capitolo~6
        (array, puntatori e stringhe).
  \item \textbf{AddressSanitizer} --- documentazione di GCC
        (\code{-fsanitize=address}). \\
        \urlx{https://gcc.gnu.org/onlinedocs/gcc/Instrumentation-Options.html}
  \item \textbf{Dynamic memory management} --- cppreference.com. \\
        \urlx{https://en.cppreference.com/w/c/memory}
\end{itemize}
